\documentclass[11pt]{article}
\usepackage[margin=1in]{geometry}
\usepackage{amsmath}
\usepackage{booktabs}
\usepackage{hyperref}
\usepackage{url}

\title{Memo 10: Concentration-Regime Settling Velocity Formulation\\
       for Fine Sediment and MP-Floc Coupling}
\author{FVCOM-MP WaterPACT Investigation}
\date{May 2026}

\begin{document}
\sloppy
\maketitle

\section*{Purpose}

This memo records a replacement direction for the current ORIG\_SED plus
Stokes-equivalent floc-size workflow.  The immediate problem is that a pure
Stokes inversion of a capped settling velocity is only a hydraulic diagnostic;
it is not a physically consistent floc-size estimate.  The alternative reviewed
here is the concentration-regime settling velocity formulation used by
Shiravani et al. (2023) for microplastic interaction with fine sediment under
estuarine conditions.  The local reference implementation is
\texttt{MPTM\_FSK\_Veloc.m}.

\section*{Sources Reviewed}

The paper is:
\begin{quote}
Shiravani, G., Oberrecht, D., Roscher, L., Kernchen, S., Halbach, M.,
Gerriets, M., Scholz-Bottcher, B.M., Gerdts, G., Badewien, T.H., and
Wurpts, A. (2023). Numerical modeling of microplastic interaction with fine
sediment under estuarine conditions. \emph{Water Research}, 231, 119564.
\url{https://doi.org/10.1016/j.watres.2022.119564}
\end{quote}

ScienceDirect lists the article as open access in \emph{Water Research}, volume
231, article 119564.  The paper highlights a new approach for MP settling
velocity and describes MP-fine sediment interaction through five settling
velocity regions.  In the abstract, the key process is described as integration
of small MP particles into fine sediment aggregates above a threshold sediment
concentration, causing the effective MP settling velocity to approach that of
sediment aggregates.

The local MATLAB script \texttt{MPTM\_FSK\_Veloc.m} implements the same
concept as a piecewise function of suspended fine sediment concentration.  This
memo follows the script constants and algebra because it is the version
available for direct implementation.

\section*{Conceptual Difference from the Present ORIG/Stokes Approach}

The current ORIG-based prototype reconstructs a local cohesive settling speed
from concentration, salinity, turbulence intensity, and primary sediment size.
Then it inverts the settling speed through Stokes law to obtain an
``equivalent'' floc diameter.  This becomes inconsistent when the reconstructed
settling speed is capped: many different physical sizes can collapse to the
same capped velocity, and Stokes inversion then returns the same small
hydraulic diameter.

The concentration-regime formulation avoids that step.  It directly prescribes
a settling speed as a function of local fine sediment concentration:
\begin{itemize}
  \item inert region at low sediment concentration;
  \item transition region where MP/sediment interaction begins;
  \item flocculation region where settling speed rises with concentration;
  \item hindered-settling region where settling speed decreases with
        concentration;
  \item consolidation region at very high concentration.
\end{itemize}

This is therefore better suited as a replacement for the effective settling
velocity signal.  It should not automatically be interpreted as a floc diameter
law.

\section*{Piecewise Formula in the MATLAB Script}

Let $C$ be fine sediment concentration in kg m$^{-3}$ and $w_s(C)$ be the
local effective settling speed in m s$^{-1}$.  The script constants are:

\begin{table}[h]
\centering
\caption{Constants extracted from \texttt{MPTM\_FSK\_Veloc.m}.}
\begin{tabular}{lll}
\toprule
Script name & Meaning in this memo & Value \\
\midrule
\texttt{V0SED} & inert MP settling speed & 1.67 m d$^{-1}$ \\
\texttt{V0SED0} & inert MP settling speed & $1.93287\times10^{-5}$ m s$^{-1}$ \\
\texttt{Reza} & lower interaction threshold $C_f$ & 0.1 kg m$^{-3}$ \\
\texttt{ch2Sedc} & transition/flocculation boundary $C_2$ & 1.0 kg m$^{-3}$ \\
\texttt{csoil} & hindered-settling threshold $C_h$ & 8.0 kg m$^{-3}$ \\
\texttt{fcgel} & gel concentration $C_{\mathrm{gel}}$ & 50.0 kg m$^{-3}$ \\
\texttt{cmax} & upper concentration limit $C_{\max}$ & 250.0 kg m$^{-3}$ \\
\texttt{ws0} & high-flocculation reference speed & $3.0\times10^{-3}$ m s$^{-1}$ \\
\texttt{fgws} & gel/consolidation reference speed & $1.0\times10^{-5}$ m s$^{-1}$ \\
\texttt{fmfloc} & flocculation exponent $m$ & 1.1 \\
\texttt{fbhis} & hindered-settling exponent $b_h$ & 2.0 \\
\texttt{cconsol} & consolidation exponent $\beta$ & 3.0 \\
\texttt{wsmin} & flocculation-zone lower bound & $1.0\times10^{-4}$ m s$^{-1}$ \\
\bottomrule
\end{tabular}
\end{table}

The script first computes
\begin{equation}
  a_h =
  \frac{1-\left(w_{\mathrm{gel}}/w_{s0}\right)^{1/b_h}}
       {C_{\mathrm{gel}}},
\end{equation}
where $w_{\mathrm{gel}}=\texttt{fgws}$ and $w_{s0}=\texttt{ws0}$.  It then
sets the flocculation coefficient
\begin{equation}
  k_1 =
  \frac{w_{s0}\left(1-a_h C_h\right)^{b_h}}
       {C_h^m}.
\end{equation}
For the provided constants, $a_h=1.88453\times10^{-2}$ and
$k_1=2.19675\times10^{-4}$ in SI units.

For the current \texttt{V0SED}=1.67 m d$^{-1}$ case, the script chooses
$C_2=1.0$ kg m$^{-3}$ and a cubic transition exponent $n=3$.  The transition
coefficient is
\begin{equation}
  m_R = \frac{k_1-w_i}{(C_2-C_f)^n},
\end{equation}
where $w_i=\texttt{V0SED0}$.

The resulting piecewise function is:
\begin{equation}
w_s(C)=
\begin{cases}
w_i, & 0 < C < C_f, \\
w_i + m_R(C-C_f)^3, & C_f \le C < C_2, \\
\max(k_1 C^m, w_{\min}), & C_2 \le C < C_h, \\
w_{s0}(1-a_h C)^{b_h}, & C_h \le C < C_{\mathrm{gel}}, \\
w_{\mathrm{gel}}\left(C/C_{\mathrm{gel}}\right)^{-\beta},
  & C_{\mathrm{gel}} \le C < C_{\max}, \\
0, & C \le 0 \text{ or } C \ge C_{\max}.
\end{cases}
\end{equation}

\section*{Representative Values}

Table~\ref{tab:values} shows selected values from the local MATLAB
implementation after applying the script constants above.

\begin{table}[h]
\centering
\caption{Representative settling speeds from \texttt{MPTM\_FSK\_Veloc.m}.}
\label{tab:values}
\begin{tabular}{rlrr}
\toprule
$C$ (kg m$^{-3}$) & Regime & $w_s$ (m s$^{-1}$) & $w_s$ (m d$^{-1}$) \\
\midrule
0.01 & inert & $1.9329\times10^{-5}$ & 1.67 \\
0.10 & transition start & $1.9329\times10^{-5}$ & 1.67 \\
0.50 & transition & $3.6917\times10^{-5}$ & 3.19 \\
1.00 & flocculation start & $2.1968\times10^{-4}$ & 18.98 \\
2.00 & flocculation & $4.7088\times10^{-4}$ & 40.68 \\
4.00 & flocculation & $1.0094\times10^{-3}$ & 87.21 \\
8.00 & hindered start & $2.1636\times10^{-3}$ & 186.94 \\
20.0 & hindered & $1.1647\times10^{-3}$ & 100.63 \\
50.0 & consolidation start & $1.0000\times10^{-5}$ & 0.864 \\
100.0 & consolidation & $1.2500\times10^{-6}$ & 0.108 \\
\bottomrule
\end{tabular}
\end{table}

This curve is qualitatively different from the current ORIG\_SED formula.  It
does not depend on salinity, local turbulence intensity $G$, or primary
sediment $D_{50}$.  It instead treats concentration regime as the controlling
variable.

\section*{Updated Modeling Scope}

The paper under development focuses on MP-floc interaction, not on resolving
floc population dynamics as the main scientific product.  Therefore, a full
CSTMS/flocmod implementation is not required for the present study if we state
the closure carefully.  A defensible pathway is to represent the observed
estuarine floc field empirically with a small number of cohesive sediment/floc
classes and to use the concentration-regime settling formulation as the
effective settling-speed closure for MP-floc interaction.

Under this scope:
\begin{itemize}
  \item the cohesive sediment classes are interpreted as empirical fine
        sediment/floc classes, not as sand grains;
  \item class properties and source strengths are tuned so that the simulated
        turbidity maximum and representative floc-size distribution are
        reasonable compared with available Delaware estuary observations;
  \item the model is used to evaluate how plausible MP-floc interaction changes
        MP transport, water-column partitioning, deposition, and bed storage;
  \item the work does not claim to mechanistically predict floc generation,
        aggregation, and fragmentation from first principles.
\end{itemize}

This keeps the scientific claim centered on MP behavior.  Sand-sized median
grain classes remain relevant for bed structure, burial, and non-cohesive
transport, but they should not be treated as MP-floc aggregate classes.

\section*{Implications for FVCOM-MP}

The formula directly gives a concentration-dependent settling velocity.  That
is the quantity we need for:
\begin{itemize}
  \item differential-settling kernels;
  \item aggregated MP settling speed;
  \item output diagnostics such as \texttt{settle\_vel\_floc\_mp1}.
\end{itemize}

However, it does not provide a physical floc diameter.  Therefore, using
Stokes inversion to recover a diameter would reintroduce the same conceptual
problem documented in Memo 09.  The next implementation should treat
$w_s(C)$ as the primary prognostic/diagnostic coupling signal.  If the MP
aggregation kernel still requires a size-exclusion diameter, that diameter
should be supplied by a separate empirical or parameterized rule, not inferred
from pure Stokes inversion of $w_s(C)$.

\section*{Implementation Plan}

\begin{enumerate}
  \item Stay within the current ORIG\_SED-compatible sediment framework for the
        next implementation stage.  Do not move to CSTMS/flocmod yet.  The
        near-term goal is an empirical multi-class cohesive setup that supports
        MP-floc sensitivity experiments without claiming full floc-population
        dynamics.
  \item Configure two or three cohesive sediment/floc classes in the sediment
        input.  These classes should represent fine primary material and larger
        empirical floc/aggregate classes.  Their \texttt{SED\_SD50},
        \texttt{SED\_WSET}, erosion/deposition thresholds, and source strengths
        should be treated as tunable empirical parameters constrained by
        Delaware estuary turbidity and floc observations.
  \item Add a new optional concentration-regime settling method in
        \texttt{mod\_sed.F}, for example
        \texttt{SED\_WSET\_METHOD = FSK\_CONC}.  The default must remain the
        current behavior so existing cases are unchanged.  The new method
        should be available for cohesive classes.
  \item Add input parameters corresponding to the MATLAB script constants:
        \texttt{SED\_FSK\_C\_F}, \texttt{SED\_FSK\_C\_2},
        \texttt{SED\_FSK\_C\_H}, \texttt{SED\_FSK\_C\_GEL},
        \texttt{SED\_FSK\_C\_MAX}, \texttt{SED\_FSK\_WS\_INERT},
        \texttt{SED\_FSK\_WS0}, \texttt{SED\_FSK\_WS\_GEL},
        \texttt{SED\_FSK\_M\_FLOC}, \texttt{SED\_FSK\_B\_HIND},
        \texttt{SED\_FSK\_BETA\_CONSOL}, \texttt{SED\_FSK\_WSMIN}, and
        \texttt{SED\_FSK\_TRANS\_N}.  The defaults should reproduce
        \texttt{MPTM\_FSK\_Veloc.m}.
  \item Implement a scalar helper, for example
        \texttt{Calc\_Wset\_FSK\_Cell(C,wset)}, and a column helper, for
        example \texttt{Calc\_Wset\_FSK(n,c,wset)}.  The helper should use only
        concentration and the configured constants.  For multi-class use, the
        concentration argument should be either total cohesive/floc
        concentration or class-specific concentration according to a documented
        switch; the preferred first implementation is total cohesive
        concentration so the ETM concentration regime controls all floc classes.
  \item In the sediment settling routine, branch the cohesive-settling path:
        retain \texttt{Calc\_Wset\_Floc} for the current ORIG method, and call
        \texttt{Calc\_Wset\_FSK} when the new method is selected.
  \item Expose the same helper to \texttt{mod\_plast.F} so the MP module can
        reconstruct the exact same local settling signal from sediment
        concentration without exchanging \texttt{Ws} arrays.  This preserves the
        earlier design principle that MP reads local sediment concentration but
        does not require a direct settling-speed exchange.
  \item Replace the current ORIG/Stokes path in MP with the FSK concentration
        path for \texttt{PLAST\_FLOC\_EFFECTIVE}.  The FSK value should be used
        directly for differential settling, aggregated MP settling, and
        \texttt{settle\_vel\_floc\_*} output.
  \item Use the configured cohesive/floc class diameter directly as the MP
        aggregation-size variable instead of pure Stokes inversion of
        \texttt{Ws}.  This is consistent with the empirical multi-class setup:
        class diameter represents the modeled floc/aggregate size, while the
        FSK formula represents the concentration-regime settling speed.
        A concentration-weighted effective class diameter can be diagnosed for
        output and plotting.
  \item Tune the cohesive/floc source terms and class parameters so that the
        model produces a plausible ETM and a plausible floc-size enhancement
        relative to available Delaware estuary measurements.  This tuning step
        should be documented as empirical calibration for MP-floc experiments.
  \item Add output diagnostics for the selected concentration regime and the
        FSK settling speed.  Also output the concentration-weighted model-used
        floc diameter and, where useful, each class contribution to the MP
        aggregation kernel.  Any hydraulic/Stokes-equivalent diameter should be
        labeled as diagnostic only.
  \item Create new paired cases only after the method is implemented: one case
        with identical multi-class cohesive sediment setup but MP-floc coupling
        off or using the previous baseline, and one or more cases with FSK
        MP-floc coupling on for selected MP types.  This keeps the sediment
        environment fixed while evaluating MP response.
\end{enumerate}

\section*{Revised Technical Position}

The size-exclusion issue should be handled empirically rather than by pure
Stokes inversion.  The planned multi-class cohesive setup supplies the
model-used floc/aggregate diameter from the configured class diameter.  The FSK
formula supplies concentration-regime settling velocity.  These two pieces are
therefore separated:
\begin{itemize}
  \item class diameter controls collision radius, number-density conversion,
        and the MP size-exclusion rule;
  \item FSK settling velocity controls differential settling and aggregated MP
        settling/output;
  \item source and class-parameter tuning controls whether the ETM produces a
        realistic empirical floc distribution.
\end{itemize}

This approach is less mechanistic than CSTMS, but it is appropriate for a paper
whose central question is how MP-floc interaction affects selected MP types
under plausible Delaware estuary floc conditions.

\section*{Implementation Update: Single-Class FSK Case D/E}

The immediate implementation will first use a single cohesive sediment class as
a formula-isolation test before returning to the broader empirical multi-class
setup.  Case D and case E will use the same sediment and plastic geometry so
that the effect of the MP effective-floc switch can be isolated:
\begin{itemize}
  \item one cohesive class with \texttt{SED\_SD50 = 0.010} mm;
  \item \texttt{SED\_WSET\_METHOD = FSK\_CONC}, where the inert settling speed
        is computed internally from \texttt{SED\_SD50} using the same
        temperature-compatible Stokes form used by the MP effective-size
        inversion;
  \item a microfibre-like MP with \texttt{PLAST\_AAXI = 0.010} mm,
        \texttt{PLAST\_BAXI = 0.002} mm, and
        \texttt{PLAST\_CAXI = 0.002} mm;
  \item case D sets \texttt{PLAST\_FLOC\_EFFECTIVE = T};
  \item case E sets \texttt{PLAST\_FLOC\_EFFECTIVE = F}.
\end{itemize}

This single-class test intentionally keeps the Stokes-equivalent inversion in
the MP effective path.  Sediment computes the FSK concentration-regime settling
velocity, and the MP module reconstructs the same FSK settling velocity from
local sediment concentration, sediment class properties, local water density,
and temperature.  The MP module then applies the existing Stokes inversion with
\texttt{PLAST\_FLOC\_RHO\_EFF} to obtain the effective floc diameter used by the
current MP-floc interaction parameterization.  This is a diagnostic
implementation test of the FSK settling pathway, not the final claim that FSK
settling uniquely determines physical floc size.

\section*{Wu et al. Minimum-Dimension Incorporation Criterion}

The MP-floc incorporation threshold will be revised to match the size
definition used by Wu et al. (2024), \emph{Flocs as vectors for microplastics
in the aquatic environment}.  In that work, MP size for the boundary curve is
the smallest relevant particle dimension: diameter for microspheres, diameter
for microfibres, and minimum Feret diameter for irregular fragments.  Therefore,
the model should not use the volume-equivalent diameter \texttt{dequ} when
testing whether an MP can be incorporated into a floc.

The code-design split is:
\begin{itemize}
  \item \texttt{dequ} remains the size used for MP settling, collision radius,
        mass, and volume-related calculations;
  \item the incorporation threshold uses
        $\min(\texttt{aaxis},\texttt{baxis},\texttt{caxis})$, which is normally
        \texttt{caxis} under the input convention;
  \item the same minimum-dimension value is used in both aggregation
        size-exclusion and disaggregation threshold calculations.
\end{itemize}

For the planned 10 $\mu$m floc test, the Wu threshold gives
$d_{\mathrm{MP,max}}\approx3.58~\mu$m.  The selected 10 by 2 by 2 $\mu$m
microfibre has a minimum dimension of 2 $\mu$m and should therefore pass the
incorporation gate, while retaining its non-spherical geometry in the rest of
the MP physics.

\end{document}
